\documentclass[11pt]{article}
\usepackage{amsmath,amssymb,geometry}
\geometry{margin=25mm}
\usepackage[UTF8]{ctex}
\begin{document}
\section*{核心模型：具有途中容量管理的空中交通流量管理（TFMP）}
来源：Bertsimas和Stock Patterson，第1.1节，PDF第3-5页，（印刷版第408-410页），方程（2）-（8）以及TFMP目标。

\paragraph{数据。}
设$F,K,J$为航班、机场和途经区域，设$\mathcal T=\{1,\ldots,H\}$为有限时间网格。航班$f$具有固定有序路径$P_f=(P(f,1),\ldots,P(f,N_f))$，起飞机场为$o_f=P(f,1)$，目的地机场为$d_f^{loc}=P(f,N_f)$。每个中间路径元素属于$J$。设$\mathcal C$包含由同一架飞机连续飞行的$(f',f)$对；它们的连接机场为$d_{f'}^{loc}=o_f$。预定的起飞/到达时间为$d_f,r_f$，地面/空中延误率为$c_f^g,c_f^a$，$s_{f'}$为$f'$之后的周转时间。正整数$l_{fj}$是从路径元素$j$到其后续元素的最小穿越时间。容量分别为$D_k(t)$用于起飞、$A_k(t)$用于到达，以及$S_j(t)$用于同时占用区域。对于每架航班/路径元素，$T_f^j=\{\underline T_f^j,\ldots,\overline T_f^j\}$为提供的可行到达时间窗口。
\paragraph{变量和边界约定。}
$w_{ft}^j\in\{0,1\}$表示在时间$t$时到达元素$j$的位置。  
使用源代码的窗口边界值。
\[
w_{ft}^j=0\quad(t<\underline T_f^j),\qquad
w_{ft}^j=1\quad(t\ge\overline T_f^j).
\]
因此，最后一个可行的时间被固定为 1，这样就能确保每趟航班都能经过每一个元素。定义这些表达式。
\[
D_f=\sum_{t\in T_f^{o_f}}t(w_{ft}^{o_f}-w_{f,t-1}^{o_f}),\quad
R_f=\sum_{t\in T_f^{d_f^{loc}}}t(w_{ft}^{d_f^{loc}}-w_{f,t-1}^{d_f^{loc}}),
\quad g_f=D_f-d_f,\quad a_f=R_f-r_f-g_f.
\]
\paragraph{一个完整的优化模型。}
\[
\min_w\ \sum_{f\in F}(c_f^g g_f+c_f^a a_f).
\]
在以下条件下：
\begin{align}
\sum_{f:o_f=k}(w_{ft}^k-w_{f,t-1}^k)&\le D_k(t)
 &&k\in K,t\in\mathcal T,\tag{2}\\
\sum_{f:d_f^{loc}=k}(w_{ft}^k-w_{f,t-1}^k)&\le A_k(t)
 &&k\in K,t\in\mathcal T,\tag{3}\\
\sum_{\substack{f,i:P(f,i)=j\\j'=P(f,i+1),\ i<N_f}}
 (w_{ft}^{j}-w_{ft}^{j'})&\le S_j(t)
 &&j\in J,t\in\mathcal T,\tag{4}\\
w_{f,t+l_{fj}}^{j'}&\le w_{ft}^{j}
 &&f\in F,\ j=P(f,i),\ j'=P(f,i+1),\ i<N_f,\ t\in T_f^j,\tag{5}\\
w_{ft}^{k}&\le w_{f',t-s_{f'}}^{k}
 &&(f',f)\in\mathcal C,\ k=o_f,\ t\in T_f^k,\tag{6}\\
w_{ft}^{j}&\ge w_{f,t-1}^{j}
 &&f\in F,j\in P_f,t\in T_f^j,\tag{7}\\
w_{ft}^{j}&\in\{0,1\}
 &&f\in F,j\in P_f,t\in T_f^j.\tag{8}
\end{align}
所有偏移的索引都采用上述边界约定。目标在代数上等同于打印出的TFMP目标，包括其时间表常数。
\end{document}
